\section{FORMAL RESTATEMENT}
\textbf{Erd\H{o}s \#94; reconstruction of a known proof, 2026-09-23.}
Let $P$ be the $n$ vertices of a strictly convex polygon. For each positive distance $r$, let $f(r)$ count unordered pairs of vertices at distance $r$. Prove $\sum_r f(r)^2=O(n^3)$. We prove the explicit, more general estimate
\[
\sum_r f(r)^2\leq\frac34n^2(n-1)
\tag{1}
\]
whenever no three points of $P$ are collinear.

\section{QUICK LITERATURE AND CONTEXT CHECK}
The source labels this problem proved and attributes the general-position result to Lefmann and Thiele (1995). The current discussion supplies the elementary isosceles-triangle argument. The following is a reconstruction of that known argument, with its counting conventions and constants checked; it is not a new resolution. Strict convexity implies the general-position hypothesis used here.

\section{OBJECTIVE AND ATTACK PLAN}
Replace each global distance multiplicity by the sum of its degrees at individual points. Cauchy--Schwarz then reduces the problem to counting isosceles triangles. Their apices lie on perpendicular bisectors, so the restriction on collinear triples gives a uniform bound.

\section{WORK}
For $x\in P$ put $d_r(x)=|\{y\in P\setminus\{x\}:|x-y|=r\}|$. Each edge has two endpoints, hence $\sum_xd_r(x)=2f(r)$. Thus
\[
4f(r)^2\leq n\sum_xd_r(x)^2.
\tag{2}
\]
Let $I$ count pairs consisting of an apex $x$ and an unordered two-point base $\{y,z\}$ with $x,y,z$ distinct and $|x-y|=|x-z|$. We count an equilateral triangle three times, once for each possible apex; this is the convention required by the identity
\[
I=\sum_{x,r}\binom{d_r(x)}2.
\]
Every other vertex contributes to exactly one degree at $x$, so $\sum_{x,r}d_r(x)=n(n-1)$. Consequently
\[
\sum_{x,r}d_r(x)^2=2I+n(n-1).
\tag{3}
\]
Fix the unordered base $\{y,z\}$. Its possible apices are points of $P$ on the perpendicular bisector of $yz$. A line contains at most two points of $P$, so there are at most two apices for this base. The endpoints themselves are not on that perpendicular bisector. Therefore
\[
I\leq2\binom n2=n(n-1).
\tag{4}
\]
Summing (2), then applying (3) and (4), proves (1).

The cubic order cannot be replaced by a smaller power uniformly over convex polygons. For a regular odd $n$-gon, there are $(n-1)/2$ distances, each realized by $n$ unordered pairs, so the sum is $n^2(n-1)/2$. For a regular even $n$-gon, the first $n/2-1$ chord lengths have multiplicity $n$ and the diameter has multiplicity $n/2$, giving $n^3/2-3n^2/4$. These computations also check the factor of two in the definition of $f$.

As a consequence, if $D$ is the number of distinct distances, another application of Cauchy--Schwarz gives
\[
D\geq\frac{\binom n2^2}{\sum_r f(r)^2}\geq\frac{n-1}{3}.
\]
This consequence is weaker than Altman's $\lfloor n/2\rfloor$ bound for convex polygons; the assertion of this problem is the stronger multiplicity control (1).

\section{VERIFICATION AND SCOPE}
All counts above are exact and the proof is complete for every finite general-position set, including $n\leq2$. No numerical experiments or unproved geometric input are used. The word ``convex'' is interpreted as vertices of a convex polygon, with no extra points on its sides; the proof does not assume that arbitrary points on the boundary of a weakly convex polygon satisfy general position.

\section{SOURCES}
\begin{itemize}
\item Problem statement and current attribution: \url{https://www.erdosproblems.com/94}.
\item H. Lefmann and T. Thiele, \emph{Point sets with distinct distances}, Combinatorica 15 (1995), 379--408: \url{https://doi.org/10.1007/BF01299744}.
\item Current discussion of the elementary argument: \url{https://www.erdosproblems.com/forum/thread/94}.
\end{itemize}

\section{CONCLUSION}
\textbf{Attempt status: solved by reconstruction of a known proof.} The explicit bound (1) proves the requested estimate; no novelty is claimed.
\textbf{Completion rate estimate: 100\%.}
